\subsection{6. Bố cục quyển báo cáo}
\indent\indent Quyển báo cáo gồm có 3 phần chính là: Phần giới thiệu, Phần nội dung và Phần kết luận, cuối quyển báo cáo là phần Tài liệu tham khảo chứa những quyển sách, bài báo được trích dẫn trong báo cáo. Cụ thể phần bố cục được chia như sau:\vspace{-0.1cm}
\begin{itemize}[label={--}, itemsep=1pt]
    \item \textbf{Phần giới thiệu:} Trình bày các vấn đề hiện tại, các nghiên cứu liên quan, mục tiêu của đề tài, đối tượng và phạm vi nghiên cứu, cùng với phương pháp nghiên cứu được áp dụng.
    \item \textbf{Phần nội dung:}\vspace{-0.1cm} 
    \begin{itemize}[label={+}, itemsep=1pt]
        \item \textbf{Chương 1 -- Mô tả bài toán:} Giới thiệu chi tiết về bài toán, bao gồm các kiến thức và framework cụ được sử dụng trong quá trình nghiên cứu đề tài.
        \item \textbf{Chương 2 -- Thiết kế và cài đặt giải pháp:} Mô tả sơ lược về pipeline của hệ thống và kiến trúc của mô hình. Đồng thời mô tả phương pháp thu thập và tiền xử lý dữ liệu.
        \item \textbf{Chương 3 -- Kiểm thử và đánh giá:} Trình bày giao diện sản phẩm và thực hiện đánh giá kết quả thực nghiệm của mô hình đề ra.
    \end{itemize}
    \item \textbf{Phần kết luận:} Tóm tắt kết quả đạt được, nêu lên những hạn chế của đề tài và đề xuất hướng phát triển trong tương lai.
    \item \textbf{Tài liệu tham khảo:} Liệt kê các nguồn tài liệu được sử dụng trong quá trình nghiên cứu và thực hiện đề tài.
\end{itemize}